\subsubsection*{Соотношение между параметрами $n$, $r$, $m$, $p$}

\begin{itemize}
    \item $r$ и $m$ обозначают одно и то же: \textbf{ранг задачи}, то есть количество пар элементов, подлежащих попарному умножению ($r = m$). Это длина векторов $A$ и $B$.
    \item $p$ --- \textbf{разрядность} одного операнда (множимое и множитель). В варианте 2 $p = 4$; допустимые значения элементов векторов от $0$ до $2^p - 1$.
    \item $n$ --- \textbf{число стадий} конвейера. В сбалансированной модели по одному разряду множителя на стадию выполняется соотношение $n = p$ (число стадий равно разрядности). Результат умножения имеет разрядность $2p$.
\end{itemize}

Итого: $m = r$ (количество пар), $n = p$ (число стадий совпадает с разрядностью), $p$ задаётся вариантом задания.
